% GENERATED FILE. Edit canonical lessons, then run npm run generate:books.
% canonical-source-hash: fnv1a64:2c9b6e495e3a0ca1
% canonical-lessons: IT-C406-accelerare, IT-C406-freno, IT-C406-frenata, IT-C406-sterzare, IT-C406-sbandare

\chapter{Alla guida}
\label{ch:it-alla-guida}
\hlchaptermodality{406}

\begin{chapteropening}
\textbf{By the end of this chapter:} \emph{I can describe driving: speeding up, the brake, braking hard, steering and skidding.}

Complete the last lesson of chapter 406: I can describe driving: speeding up, the brake, braking hard, steering and skidding.
\end{chapteropening}

\section[accelerare]{accelerare — to speed up, to accelerate}

\label{lesson:IT-C406-accelerare}



\begin{quote}
Before the new one: say the Italian for the junction, then the Italian for to overtake.
\end{quote}

\subsection*{You'll want to know: accelerare}
\textbf{accelerare} — \textquotedblleft{}to speed up, to accelerate\textquotedblright{}.

\begin{quote}
Non accelerare, c'è la nebbia.
\end{quote}

\begin{quote}
\emph{Do not speed up, it is foggy.}
\end{quote}

\emph{L'acceleratore} is the accelerator pedal.

\begin{cousinweb}
\emph{accelerare} is a learned borrowing from Latin \emph{accelerāre} 'to hasten' (EtymDB), built on \emph{celer} 'swift'.
\end{cousinweb}

\begin{grammarlens}[title={What you've built}]
One more word for telling, explaining and describing at B1.
\end{grammarlens}

\subsection*{Your turn}
\begin{itemize}
\raggedright
  \item \emph{Say it:} \emph{accelerare}
  \item \emph{Say it:} \emph{accelerare}, once more
  \item \emph{Say it:} a sentence of your own with \emph{accelerare}
  \item \emph{Recall:} say the Italian for the junction, then the Italian for to overtake, then say \emph{accelerare} again
\end{itemize}

\begin{tcolorbox}[breakable,colback=teal!4,colframe=teal!35!black,title={Before you move on}]
What does accelerare mean? (\textquotedblleft{}To speed up, to accelerate\textquotedblright{}.) Say it once more.
\end{tcolorbox}
\section[il freno]{il freno — the brake}

\label{lesson:IT-C406-freno}



\begin{quote}
Before the new one: say the Italian for to overtake, then the Italian for to speed up.
\end{quote}

\subsection*{You'll want to know: il freno}
\textbf{il freno} — \textquotedblleft{}the brake\textquotedblright{}.

\begin{quote}
Il freno della bici non funziona.
\end{quote}

\begin{quote}
\emph{The brake on the bike does not work.}
\end{quote}

\emph{Il freno a mano} is the handbrake. Figuratively, \emph{senza freni} means without restraint.

\begin{cousinweb}
\emph{freno} is inherited from Latin \emph{frēnum} 'bridle, bit' (EtymDB): first the thing that held back a horse.
\end{cousinweb}

\begin{grammarlens}[title={What you've built}]
One more word for telling, explaining and describing at B1.
\end{grammarlens}

\subsection*{Your turn}
\begin{itemize}
\raggedright
  \item \emph{Say it:} \emph{il freno}
  \item \emph{Say it:} \emph{il freno}, once more
  \item \emph{Say it:} a sentence of your own with \emph{il freno}
  \item \emph{Recall:} say the Italian for to overtake, then the Italian for to speed up, then say \emph{il freno} again
\end{itemize}

\begin{tcolorbox}[breakable,colback=teal!4,colframe=teal!35!black,title={Before you move on}]
What does il freno mean? (\textquotedblleft{}The brake\textquotedblright{}.) Say it once more.
\end{tcolorbox}
\section[la frenata]{la frenata — the braking; a sudden stop}

\label{lesson:IT-C406-frenata}



\begin{quote}
Before the new one: say the Italian for to speed up, then the Italian for the brake.
\end{quote}

\subsection*{You'll want to know: la frenata}
\textbf{la frenata} — \textquotedblleft{}the braking; a sudden stop\textquotedblright{}.

\begin{quote}
Ho fatto una frenata e il gatto è scappato.
\end{quote}

\begin{quote}
\emph{I braked hard and the cat ran off.}
\end{quote}

\emph{Frenare} is to brake; \emph{una frenata brusca} is slamming on the brakes.

\begin{cousinweb}
\emph{frenata} is built on \emph{frenare} 'to brake', from \emph{freno} 'brake', which is inherited from Latin \emph{frēnum} 'bridle' (EtymDB).
\end{cousinweb}

\begin{grammarlens}[title={What you've built}]
One more word for telling, explaining and describing at B1.
\end{grammarlens}

\subsection*{Your turn}
\begin{itemize}
\raggedright
  \item \emph{Say it:} \emph{la frenata}
  \item \emph{Say it:} \emph{la frenata}, once more
  \item \emph{Say it:} a sentence of your own with \emph{la frenata}
  \item \emph{Recall:} say the Italian for to speed up, then the Italian for the brake, then say \emph{la frenata} again
\end{itemize}

\begin{tcolorbox}[breakable,colback=teal!4,colframe=teal!35!black,title={Before you move on}]
What does la frenata mean? (\textquotedblleft{}The braking; a sudden stop\textquotedblright{}.) Say it once more.
\end{tcolorbox}
\section[sterzare]{sterzare — to steer, to swerve}

\label{lesson:IT-C406-sterzare}



\begin{quote}
Before the new one: say the Italian for the brake, then the Italian for the braking.
\end{quote}

\subsection*{You'll want to know: sterzare}
\textbf{sterzare} — \textquotedblleft{}to steer, to swerve\textquotedblright{}.

\begin{quote}
Ho sterzato per evitare un cane.
\end{quote}

\begin{quote}
\emph{I swerved to avoid a dog.}
\end{quote}

\emph{Lo sterzo} is the steering; \emph{sterzare a sinistra} — to turn the wheel left.

\begin{cousinweb}
\emph{sterzare} is built on \emph{sterzo} 'steering', which was borrowed from Lombardic \emph{sterz} 'the handle of a plough' (EtymDB): the farmer guided the plough by it.
\end{cousinweb}

\begin{grammarlens}[title={What you've built}]
One more word for telling, explaining and describing at B1.
\end{grammarlens}

\subsection*{Your turn}
\begin{itemize}
\raggedright
  \item \emph{Say it:} \emph{sterzare}
  \item \emph{Say it:} \emph{sterzare}, once more
  \item \emph{Say it:} a sentence of your own with \emph{sterzare}
  \item \emph{Recall:} say the Italian for the brake, then the Italian for the braking, then say \emph{sterzare} again
\end{itemize}

\begin{tcolorbox}[breakable,colback=teal!4,colframe=teal!35!black,title={Before you move on}]
What does sterzare mean? (\textquotedblleft{}To steer, to swerve\textquotedblright{}.) Say it once more.
\end{tcolorbox}
\section[sbandare]{sbandare — to skid, to swerve off course}

\label{lesson:IT-C406-sbandare}



\begin{quote}
Before the new one: say the Italian for the braking, then the Italian for to steer.
\end{quote}

\subsection*{You'll want to know: sbandare}
\textbf{sbandare} — \textquotedblleft{}to skid, to swerve off course\textquotedblright{}.

\begin{quote}
Con il ghiaccio la macchina ha sbandato.
\end{quote}

\begin{quote}
\emph{On the ice the car skidded.}
\end{quote}

Also figurative: \emph{un ragazzo sbandato} is a young man who has lost his way.

\begin{cousinweb}
\emph{sbandare} is built from \emph{s-} and \emph{banda} 'side': going off to one side. The origin of \emph{banda} in this sense is not settled; it is usually linked to a Germanic word that reached Italian through Old Provençal.
\end{cousinweb}

\begin{grammarlens}[title={What you've built}]
One more word for telling, explaining and describing at B1.
\end{grammarlens}

\subsection*{Your turn}
\begin{itemize}
\raggedright
  \item \emph{Say it:} \emph{sbandare}
  \item \emph{Say it:} \emph{sbandare}, once more
  \item \emph{Say it:} a sentence of your own with \emph{sbandare}
  \item \emph{Recall:} say the Italian for the braking, then the Italian for to steer, then say \emph{sbandare} again
\end{itemize}

\begin{tcolorbox}[breakable,colback=teal!4,colframe=teal!35!black,title={Before you move on}]
What does sbandare mean? (\textquotedblleft{}To skid, to swerve off course\textquotedblright{}.) Say it once more.
\end{tcolorbox}
